\subsection{Deep vs. Shallow Structural Cloning}

Copying a list is not the same operation as assigning another name to the same list. This distinction is essential because a list is a mutable object. If two names point to the same list, a mutation through one name is visible through the other.

The issue becomes even more important with nested lists. A copy may duplicate only the outer list structure while still sharing the inner list objects.

\begin{verbatim}
a = ["spam", "eggs"]

b = a

print(f"a: {a!r}")
print(f"b: {b!r}")

print()
print(f"id(a): {id(a)}")
print(f"id(b): {id(b)}")
print(f"a is b: {a is b}")
\end{verbatim}

Possible output:

\begin{verbatim}
a: ['spam', 'eggs']
b: ['spam', 'eggs']

id(a): 140619786870656
id(b): 140619786870656
a is b: True
\end{verbatim}

There are two names, but only one list object.

\subsubsection{Assignment Is Not Copying: Shared List References with \texttt{b = a}}

The assignment \texttt{b = a} does not create a new list. It binds the name \texttt{b} to the same object already referenced by \texttt{a}.

\begin{verbatim}
a = ["Homer", "Marge"]

b = a

b.append("Bart")

print(f"a: {a!r}")
print(f"b: {b!r}")

print()
print(f"a is b: {a is b}")
\end{verbatim}

Possible output:

\begin{verbatim}
a: ['Homer', 'Marge', 'Bart']
b: ['Homer', 'Marge', 'Bart']

a is b: True
\end{verbatim}

The mutation was performed through \texttt{b}, but \texttt{a} sees it because both names refer to the same list object.

This is not special to \texttt{append()}. Any in-place mutation has the same visibility.

\begin{verbatim}
a = ["Jerry", "Elaine", "George"]

b = a

b[2] = "Kramer"

print(f"a: {a!r}")
print(f"b: {b!r}")
print(f"a is b: {a is b}")
\end{verbatim}

Possible output:

\begin{verbatim}
a: ['Jerry', 'Elaine', 'Kramer']
b: ['Jerry', 'Elaine', 'Kramer']
a is b: True
\end{verbatim}

The list slot at index \texttt{2} was replaced in the one shared list object.

The same identity relation can be drawn as:

\begin{verbatim}
a  ----+
       |
       v
    list object: ["Jerry", "Elaine", "Kramer"]
       ^
       |
b  ----+
\end{verbatim}

The practical rule is:

\begin{quote}
Assignment creates another reference to the same object. It does not clone the object.
\end{quote}

\subsubsection{Shallow Copies: \texttt{a[:]}, \texttt{list(a)}, and \texttt{.copy()}}

A shallow copy creates a new outer list. The new list has its own slot array, but the slots refer to the same element objects as the original list.

There are several common ways to create a shallow list copy:

\begin{itemize}
    \item \texttt{a[:]};
    \item \texttt{list(a)};
    \item \texttt{a.copy()}.
\end{itemize}

\begin{verbatim}
a = ["spam", "eggs", 42]

b = a[:]
c = list(a)
d = a.copy()

print(f"a: {a!r}")
print(f"b: {b!r}")
print(f"c: {c!r}")
print(f"d: {d!r}")

print()
print(f"a is b: {a is b}")
print(f"a is c: {a is c}")
print(f"a is d: {a is d}")
\end{verbatim}

Possible output:

\begin{verbatim}
a: ['spam', 'eggs', 42]
b: ['spam', 'eggs', 42]
c: ['spam', 'eggs', 42]
d: ['spam', 'eggs', 42]

a is b: False
a is c: False
a is d: False
\end{verbatim}

The outer lists are different objects.

A mutation of the outer structure of one list does not change the outer structure of the other list.

\begin{verbatim}
a = ["spam", "eggs", 42]

b = a.copy()

b.append("bacon")
b[0] = "beans"

print(f"a: {a!r}")
print(f"b: {b!r}")

print()
print(f"a is b: {a is b}")
\end{verbatim}

Possible output:

\begin{verbatim}
a: ['spam', 'eggs', 42]
b: ['beans', 'eggs', 42, 'bacon']

a is b: False
\end{verbatim}

The shallow copy protected the outer list structure. The name \texttt{b} refers to a separate list.

However, shallow copying does not copy the element objects themselves. It copies the references.

\begin{verbatim}
a = ["spam", "eggs", 42]
b = a.copy()

print("outer list identities")
print(f"id(a): {id(a)}")
print(f"id(b): {id(b)}")

print()
print("element identities")
for index in range(len(a)):
    print(f"index={index} id(a[index])={id(a[index])} "
          f"id(b[index])={id(b[index])} same={a[index] is b[index]}")
\end{verbatim}

Possible output:

\begin{verbatim}
outer list identities
id(a): 140619786870656
id(b): 140619786871040

element identities
index=0 id(a[index])=140619786560368 id(b[index])=140619786560368 same=True
index=1 id(a[index])=140619786560432 id(b[index])=140619786560432 same=True
index=2 id(a[index])=140619796012304 id(b[index])=140619796012304 same=True
\end{verbatim}

The outer list is new. The element references are shared.

For immutable elements such as strings, integers, floats, and tuples of immutable objects, this sharing is usually harmless because the shared element objects cannot be changed in place.

\begin{verbatim}
a = ["Ni", 42]
b = a.copy()

b[0] = "Ekki-ekki"
b[1] = 100

print(f"a: {a!r}")
print(f"b: {b!r}")
\end{verbatim}

Possible output:

\begin{verbatim}
a: ['Ni', 42]
b: ['Ekki-ekki', 100]
\end{verbatim}

The list slots were replaced in \texttt{b}. The original list \texttt{a} kept its own slots.

\subsubsection{Nested Structures: Why Shallow Copying Fails for Lists Inside Lists}

A shallow copy becomes dangerous when the list contains mutable objects, especially other lists.

\begin{verbatim}
a = [
    ["Homer", "Marge"],
    ["Bart", "Lisa"],
]

b = a.copy()

print("before mutation")
print(f"a: {a!r}")
print(f"b: {b!r}")

print()
print(f"a is b: {a is b}")
print(f"a[0] is b[0]: {a[0] is b[0]}")
print(f"a[1] is b[1]: {a[1] is b[1]}")
\end{verbatim}

Possible output:

\begin{verbatim}
before mutation
a: [['Homer', 'Marge'], ['Bart', 'Lisa']]
b: [['Homer', 'Marge'], ['Bart', 'Lisa']]

a is b: False
a[0] is b[0]: True
a[1] is b[1]: True
\end{verbatim}

The outer list was copied. The inner lists were not copied.

A mutation inside an inner list is therefore visible through both outer lists.

\begin{verbatim}
a = [
    ["Homer", "Marge"],
    ["Bart", "Lisa"],
]

b = a.copy()

b[0].append("Maggie")

print(f"a: {a!r}")
print(f"b: {b!r}")

print()
print(f"a[0] is b[0]: {a[0] is b[0]}")
\end{verbatim}

Possible output:

\begin{verbatim}
a: [['Homer', 'Marge', 'Maggie'], ['Bart', 'Lisa']]
b: [['Homer', 'Marge', 'Maggie'], ['Bart', 'Lisa']]

a[0] is b[0]: True
\end{verbatim}

This is the common shallow-copy surprise. The outer container is separate, but the nested mutable objects are still shared.

The structure is:

\begin{verbatim}
a  ---> outer list A
          slot 0 ----+
                     |
                     v
                  inner list ["Homer", "Marge", "Maggie"]
                     ^
                     |
b  ---> outer list B  |
          slot 0 ----+
\end{verbatim}

Replacing an inner list is different from mutating the shared inner list.

\begin{verbatim}
a = [
    ["Homer", "Marge"],
    ["Bart", "Lisa"],
]

b = a.copy()

b[0] = ["Ned", "Maude"]

print(f"a: {a!r}")
print(f"b: {b!r}")

print()
print(f"a[0] is b[0]: {a[0] is b[0]}")
print(f"a[1] is b[1]: {a[1] is b[1]}")
\end{verbatim}

Possible output:

\begin{verbatim}
a: [['Homer', 'Marge'], ['Bart', 'Lisa']]
b: [['Ned', 'Maude'], ['Bart', 'Lisa']]

a[0] is b[0]: False
a[1] is b[1]: True
\end{verbatim}

The first inner-list slot in \texttt{b} was redirected to a new list. The second inner list is still shared.

This is the key distinction:

\begin{itemize}
    \item replacing a slot changes one outer list;
    \item mutating a shared inner object is visible through all references to that object.
\end{itemize}

A repeated-list construction can create the same problem even without copying.

\begin{verbatim}
row = [0, 0, 0]

grid = [row] * 3

print("before")
print(grid)

grid[0][1] = 42

print()
print("after")
print(grid)

print()
for index, item in enumerate(grid):
    print(f"row {index}: id={id(item)} value={item}")
\end{verbatim}

Possible output:

\begin{verbatim}
before
[[0, 0, 0], [0, 0, 0], [0, 0, 0]]

after
[[0, 42, 0], [0, 42, 0], [0, 42, 0]]

row 0: id=140619786871232 value=[0, 42, 0]
row 1: id=140619786871232 value=[0, 42, 0]
row 2: id=140619786871232 value=[0, 42, 0]
\end{verbatim}

The outer list contains three references to the same row object.

A safer construction creates a new row each time.

\begin{verbatim}
grid = []

for _ in range(3):
    grid.append([0, 0, 0])

grid[0][1] = 42

print(grid)

print()
for index, item in enumerate(grid):
    print(f"row {index}: id={id(item)} value={item}")
\end{verbatim}

Possible output:

\begin{verbatim}
[[0, 42, 0], [0, 0, 0], [0, 0, 0]]

row 0: id=140619786871232 value=[0, 42, 0]
row 1: id=140619786871296 value=[0, 0, 0]
row 2: id=140619786871360 value=[0, 0, 0]
\end{verbatim}

Each row is now a separate list object.

The practical rule is:

\begin{quote}
A shallow copy is enough for a flat list whose elements will not be mutated in place. It is not enough when the list contains mutable nested structures that must become independent.
\end{quote}

\subsubsection{Deep Copying with \texttt{copy.deepcopy()}}

The standard library module \texttt{copy} provides \texttt{deepcopy()}. A deep copy recursively copies the outer object and the objects contained inside it.

\begin{verbatim}
import copy

a = [
    ["Homer", "Marge"],
    ["Bart", "Lisa"],
]

b = copy.deepcopy(a)

print("before mutation")
print(f"a: {a!r}")
print(f"b: {b!r}")

print()
print(f"a is b: {a is b}")
print(f"a[0] is b[0]: {a[0] is b[0]}")
print(f"a[1] is b[1]: {a[1] is b[1]}")
\end{verbatim}

Possible output:

\begin{verbatim}
before mutation
a: [['Homer', 'Marge'], ['Bart', 'Lisa']]
b: [['Homer', 'Marge'], ['Bart', 'Lisa']]

a is b: False
a[0] is b[0]: False
a[1] is b[1]: False
\end{verbatim}

The outer list is different, and the inner lists are also different.

Now an inner mutation remains local to the copied structure.

\begin{verbatim}
import copy

a = [
    ["Homer", "Marge"],
    ["Bart", "Lisa"],
]

b = copy.deepcopy(a)

b[0].append("Maggie")

print(f"a: {a!r}")
print(f"b: {b!r}")

print()
print(f"a[0] is b[0]: {a[0] is b[0]}")
\end{verbatim}

Possible output:

\begin{verbatim}
a: [['Homer', 'Marge'], ['Bart', 'Lisa']]
b: [['Homer', 'Marge', 'Maggie'], ['Bart', 'Lisa']]

a[0] is b[0]: False
\end{verbatim}

This is the behavior wanted when nested mutable structures must be independent.

The difference between assignment, shallow copy, and deep copy can be shown together.

\begin{verbatim}
import copy

a = [["spam"], ["eggs"]]

b = a
c = a.copy()
d = copy.deepcopy(a)

a[0].append(42)

print(f"a: {a!r}")
print(f"b: {b!r}")
print(f"c: {c!r}")
print(f"d: {d!r}")

print()
print(f"b is a:       {b is a}")
print(f"c is a:       {c is a}")
print(f"d is a:       {d is a}")

print()
print(f"c[0] is a[0]: {c[0] is a[0]}")
print(f"d[0] is a[0]: {d[0] is a[0]}")
\end{verbatim}

Possible output:

\begin{verbatim}
a: [['spam', 42], ['eggs']]
b: [['spam', 42], ['eggs']]
c: [['spam', 42], ['eggs']]
d: [['spam'], ['eggs']]

b is a:       True
c is a:       False
d is a:       False

c[0] is a[0]: True
d[0] is a[0]: False
\end{verbatim}

The result is:

\begin{itemize}
    \item \texttt{b = a}: same outer list;
    \item \texttt{c = a.copy()}: different outer list, shared inner lists;
    \item \texttt{d = copy.deepcopy(a)}: different outer list, different inner lists.
\end{itemize}

Deep copying should not be used blindly. It may be unnecessary for flat structures, and it can be expensive for large nested object graphs. But when independent nested mutable structures are required, it is the standard tool.

The practical rule is:

\begin{quote}
Use assignment when another name should refer to the same list. Use a shallow copy when only the outer list must be independent. Use \texttt{copy.deepcopy()} when nested mutable objects must also be independent.
\end{quote}